%% =====================================================================
%%  core/config/banda-prensa.tex  --  Banda institucional de prensa
%% ---------------------------------------------------------------------
%%  Pie de pagina a SANGRE COMPLETA para los documentos publicos
%%  (comunicados, pronunciamientos, notas de prensa, anuncios), en el
%%  espiritu de las piezas de despacho parlamentario: franja textil
%%  "pallay" + banda blanca con la identidad del despacho + retrato
%%  oficial que sobresale de la banda.
%%
%%  Anatomia (de arriba hacia abajo):
%%     franja pallay   campo puka con rombos blanco/q'illu alternados
%%     linea q'illu    filete fino de transicion
%%     banda blanca    RETRATO (izq., sobresale) | nombre + cajita
%%                     distrito + contacto | consigna a color (der.)
%%
%%  El RETRATO es la foto oficial de cuerpo recortada con FONDO
%%  TRANSPARENTE (la fija el perfil con \setbandafoto): se apoya en el borde
%%  inferior de la pagina y sobresale por encima de la franja, ocupando
%%  el lugar que antes tenia el logo de la Camara (pedido de Edison,
%%  2026-08-06). El logo institucional queda en el membrete superior.
%%
%%  Se dibuja como overlay TikZ anclado a la pagina fisica (sangre
%%  completa, ignora margenes) desde el pie fancyhdr: solo las paginas
%%  con la banda activa la llevan. `remember picture` exige DOS pasadas
%%  de compilacion -- las mismas que build.sh ya ejecuta (CFG_PASADAS).
%%
%%  Comandos del documento:
%%     \conbandaprensa              activa la banda (pie normal fuera)
%%     \sinbandaprensa              restaura el pie institucional
%%     \sinfotobanda / \confotobanda       retrato de la banda
%%     \sinconsignabanda / \conconsignabanda   arte de la consigna
%%     \setbandaaltura{2.90cm}      alto de la banda blanca
%%     \setbandafotoaltura{4.15cm}  alto del retrato (> banda: sobresale)
%%     \setbandafoto{archivo}       otro retrato (PNG transparente)
%%     \setbandaconsigna{archivo}   otro arte de consigna
%%
%%  OJO: la banda REEMPLAZA al pie institucional; el contacto del
%%  despacho (correo, celular) viaja dentro de la banda misma.
%% =====================================================================

\RequirePackage{tikz}

%% --- Ajustes finos ----------------------------------------------------
\newlength{\despachoBandaAlto}
\setlength{\despachoBandaAlto}{2.90cm}
\newcommand{\setbandaaltura}[1]{\setlength{\despachoBandaAlto}{#1}}

\newlength{\despachoBandaFotoAlto}
\setlength{\despachoBandaFotoAlto}{4.50cm}
\newcommand{\setbandafotoaltura}[1]{\setlength{\despachoBandaFotoAlto}{#1}}

\newcommand{\despachoBandaFoto}{retrato-banda.png}   % lo fija el perfil (\setbandafoto)
\newcommand{\setbandafoto}[1]{\renewcommand{\despachoBandaFoto}{#1}}

\newcommand{\despachoBandaConsigna}{consigna.png}   % lo fija el perfil (\setbandaconsigna)
\newcommand{\setbandaconsigna}[1]{\renewcommand{\despachoBandaConsigna}{#1}}

\newif\ifdespacho@bandafoto
\despacho@bandafototrue
\newcommand{\sinfotobanda}{\despacho@bandafotofalse}
\newcommand{\confotobanda}{\despacho@bandafototrue}

\newif\ifdespacho@bandaconsigna
\despacho@bandaconsignatrue
\newcommand{\sinconsignabanda}{\despacho@bandaconsignafalse}
\newcommand{\conconsignabanda}{\despacho@bandaconsignatrue}

%% --- La banda ---------------------------------------------------------
%%  Overlay de dimension CERO (tikzpicture overlay): el pie fancyhdr no
%%  crece; todo se pinta anclado al borde fisico inferior de la pagina.
%%  Orden de dibujado: banda -> franja -> texto/consigna -> RETRATO
%%  encima (asi la figura recortada se superpone limpiamente a la franja
%%  pallay al romperla por la izquierda).
\newcommand{\bandaprensa}{%
  % La columna de texto arranca despues del retrato si lo hay; si el
  % documento pidio \sinfotobanda, se corre al margen izquierdo.
  \ifdespacho@bandafoto
    \def\despachoBandaTextoX{4.70}%
  \else
    \def\despachoBandaTextoX{1.40}%
  \fi
  \begin{tikzpicture}[remember picture, overlay]
    \begin{scope}[shift={(current page.south west)}]
      % --- Banda blanca (tapa la marca de agua en esta zona) ---------
      \fill[white] (0,0) rectangle (\paperwidth,\despachoBandaAlto);
      % --- Remate superior: linea q'illu + franja pallay -------------
      \begin{scope}[shift={(0,\despachoBandaAlto)}]
        \fill[colorbandalinea] (0,0)    rectangle (\paperwidth,0.09);
        \fill[colorpallay]     (0,0.09) rectangle (\paperwidth,0.45);
        % rombos alternados blanco / q'illu, centrados en el campo
        \foreach \x in {0.30,1.50,...,21.00}{%
          \fill[white] (\x,0.175) -- (\x+0.095,0.27) -- (\x,0.365)
                       -- (\x-0.095,0.27) -- cycle;}
        \foreach \x in {0.90,2.10,...,21.00}{%
          \fill[colorpallayacento] (\x,0.175) -- (\x+0.095,0.27)
                       -- (\x,0.365) -- (\x-0.095,0.27) -- cycle;}
      \end{scope}
      % --- Bloque de identidad + contacto ----------------------------
      \node[anchor=west, inner sep=0pt] at (\despachoBandaTextoX,1.48) {%
        \begin{minipage}{9.7cm}
          \setlength{\parskip}{0pt}\setlength{\parindent}{0pt}%
          {\fuentedatossemi\fontsize{6.8}{8}\selectfont\color{neutral700}%
            \addfontfeatures{LetterSpace=14.0}%
            CONGRESO DE LA REPÚBLICA · CÁMARA DE DIPUTADOS\par}
          \vspace{4.5pt}
          {\fuentedisplaynegra\fontsize{13.5}{15}\selectfont\color{tinta}%
            \MakeUppercase{\ifmetadef{autorbanda}{\meta{autorbanda}}{\meta{autor}}}\par}
          \vspace{5pt}
          {\setlength{\fboxsep}{2.8pt}%
           \colorbox{colorbandatag}{%
             \fuentedatossemi\fontsize{7.2}{7.2}\selectfont\color{white}%
             \addfontfeatures{LetterSpace=12.0}\,%
             \ifmetadef{distrito}%
               {DIPUTADO POR \MakeUppercase{\meta{distrito}}}%
               {\MakeUppercase{\meta{cargo}}}\,}}%
          \par\vspace{5pt}
          {\fuentedatos\fontsize{6.6}{8}\selectfont\color{neutral700}%
            \ifmetadef{correo}{\meta{correo}}{}%
            \ifmetadef{celular}{\ \,\textperiodcentered\ \,Cel.~\meta{celular}}{}%
            \par}
        \end{minipage}};
      % --- Consigna a color (arte del despacho), al extremo derecho ---
      \ifdespacho@bandaconsigna
        \node[anchor=east, inner sep=0pt] at (19.90,1.48)
          {\includegraphics[height=1.15cm]{\despachoBandaConsigna}};
      \fi
      % --- Retrato oficial (fondo transparente): apoyado en el borde
      %     inferior, sobresale la banda; ocupa el lugar del ex-logo ---
      \ifdespacho@bandafoto
        \node[anchor=south west, inner sep=0pt] at (1.25,0)
          {\includegraphics[height=\despachoBandaFotoAlto]{\despachoBandaFoto}};
      \fi
    \end{scope}
  \end{tikzpicture}%
}

%% --- Activacion (mismo mecanismo que \sinpieinstitucional) -----------
\newcommand{\conbandaprensa}{\fancyfoot[C]{\bandaprensa}}
\newcommand{\sinbandaprensa}{\fancyfoot[C]{\pieinstitucional}}
